\input{generated/results_macros.tex}

\begin{abstract}
Response-aware test-set selection can couple a benchmark-construction objective to a downstream evaluation metric, but the resulting change in a headline metric depends on the selected chemical population. We combined a molecular endpoint-permutation experiment with an exact-size paired audit of six public QSAR/QSPR datasets. Real molecular records, Bemis--Murcko scaffold geometry, endpoint marginals, 20 partition seeds, and nested target-blind candidate pools were preserved; only structure--endpoint association was destroyed in 200 permutations per dataset. Under the molecular null, response-aware selection reduced the squared train--test mean-gap contribution to regression MSE in every permutation, whereas central 95\% null intervals for total RMSE effects spanned zero. For binary endpoints, the corresponding prevalence-gap component of Brier score was reduced, constant-score ROC--AUC remained invariant, and total Brier and log-loss effects were small and dataset-dependent. In an explicitly exploratory matched-budget bridge, observed mean-only RMSE effects exceeded all 200 null permutations for ESOL and FreeSolv under the shared-acyclic convention and for ESOL under singleton semantics, but not for singleton FreeSolv. The empirical learned-model audit found no corrected classification advantage, six of six primary regression advantages under the shared-acyclic convention, strong attenuation under singleton semantics, and representation-sensitive classification point estimates. Split-objective--metric coupling is therefore deterministic at the aligned metric-component level but contingent at the headline-metric level. We provide an auditable reporting checklist for response-aware QSAR benchmark construction.
\end{abstract}

\noindent\textbf{Keywords:} QSAR validation; scaffold splitting; benchmark construction; response permutation; metric coupling

\section{Introduction}

A QSAR benchmark is a measurement design rather than a neutral container for model scores. The validity of a held-out result depends not only on the predictor but also on molecular identity, train--test allocation, applicability domain, endpoint distribution, and the relation between the validation sample and the intended use population \citep{tropsha2003earnest,gramatica2007principles,netzeva2005domain,sahigara2012domain}. Recent chemometric work likewise emphasizes resampling robustness, leakage-aware scaling, and explicit validation rules \citep{lasfar2024robustness,kiraly2025leakage,camacho2026validation}. Within \textit{SAR and QSAR in Environmental Research}, prior work has shown that the choice and exchange of training and test sets can change model assessment, that cross-validation strategy can alter apparent performance, and that validation methods should be chosen together with applicability-domain analysis \citep{racz2015consistency,chatterjee2022crossvalidation,heberger2023validation}.

Modern molecular-machine-learning studies commonly use scaffold-disjoint validation to reduce direct framework overlap \citep{wu2018moleculenet,bemis1996properties}. Yet ``scaffold split'' is not a complete protocol. The same molecular universe can yield many scaffold-disjoint partitions with different test sizes, scaffold concentrations, endpoint prevalence or mean, and withheld chemical series. Comparative work has documented substantial variation across split strategies and random seeds \citep{yang2019learned,deng2023systematic}. DataSAIL and SIMPD address leakage-aware and prospective-proxy questions, respectively, while recent analyses emphasize that dataset partitioning and molecular-data definition are integral to a QSAR claim \citep{joeres2025datasail,landrum2023simpd,li2026partition,nael2026dataset}.

A specific risk arises when endpoint information enters partition selection. Suppose a target-blind pool of scaffold-disjoint candidates is generated and a test partition is then chosen because its endpoint mean is closer to the training mean. Even when test cardinality is held exactly fixed, the selection objective can align with the evaluation metric. For a constant regression predictor equal to the training mean,
\begin{equation}
\mathrm{MSE}_{\mathrm{mean}}=\mathrm{Var}(y_{\mathrm{test}})+(\bar y_{\mathrm{test}}-\bar y_{\mathrm{train}})^2.
\label{eq:mse-decomposition}
\end{equation}
Minimizing the absolute train--test mean gap necessarily reduces the second term relative to a paired candidate whenever the selected partition changes. It does not necessarily reduce total MSE because the selected test variance can also change. For a binary endpoint, a constant training-prevalence prediction $p$ evaluated on test prevalence $q$ has
\begin{equation}
\mathrm{Brier}=q(1-q)+(q-p)^2,
\label{eq:brier-decomposition}
\end{equation}
whereas constant-score ROC--AUC is 0.5 whenever both classes are represented. The relevant question is therefore not simply whether response-aware selection changes a score, but which metric component is coupled by construction and which collateral changes determine the final score.

Response permutation, often called $y$-randomization in QSAR validation, provides a natural way to isolate this mechanism while retaining the original molecular geometry \citep{rucker2007yrandomization}. Here the permutations are not used to validate a molecular predictor. Molecular structures and scaffold groups are retained, candidate split pools are unchanged, and no molecular predictor is fitted. The experiment asks what split selection alone can do when structure carries no endpoint information.

The study has four linked contributions. First, it formalizes split-objective--metric coupling for regression MSE and binary probability metrics. Second, it performs a pre-specified molecular endpoint-permutation experiment over six public datasets, 20 partition seeds, nested search budgets, and two acyclic-scaffold conventions. Third, it connects the null experiment to an existing exact-size paired empirical audit containing classical molecular predictors, a mean-only control, scaffold-semantic sensitivity, disconnected-component sensitivity, and partition-level inference. Fourth, it converts the combined evidence into a minimum-reporting checklist for response-aware QSAR benchmarks. The intervention is treated as an audit perturbation, not as a recommended universal splitter.
